%http://www.daniel-brettschneider.de/allgemein/latex-vorlage-fur-hausarbeiten-oder-abschlussarbeiten
\documentclass[12pt,a4paper,bibliography=totocnumbered,listof=totocnumbered]{scrartcl}
\usepackage[ngerman]{babel}
\usepackage[autostyle=true,german=quotes]{csquotes}
\usepackage[utf8]{inputenc}
\usepackage{amsmath}
\usepackage{amsfonts}
\usepackage{amssymb}
\usepackage{graphicx}
\usepackage{fancyhdr}
\usepackage{tabularx}
\usepackage{geometry}
\usepackage{setspace}
\usepackage[right]{eurosym}
\usepackage[printonlyused]{acronym}
\usepackage{subfig}
\usepackage{floatflt}
\usepackage[usenames,dvipsnames]{color}
\usepackage{colortbl}
\usepackage{paralist}
\usepackage{array}
\usepackage{titlesec}
\usepackage{parskip}
\usepackage[right]{eurosym}
\usepackage[subfigure,titles]{tocloft}
\usepackage[pdfpagelabels=true]{hyperref}
\usepackage[ngerman]{babel}
\usepackage{booktabs}
\usepackage{listings}
\lstset{basicstyle=\footnotesize, captionpos=b, breaklines=true, showstringspaces=false, tabsize=2, frame=lines, numbers=left, numberstyle=\tiny, xleftmargin=2em, framexleftmargin=2em}
\makeatletter
\def\l@lstlisting#1#2{\@dottedtocline{1}{0em}{1em}{\hspace{1,5em} Lst. #1}{#2}}
\makeatother
\geometry{a4paper, top=27mm, left=30mm, right=20mm, bottom=35mm, headsep=10mm, footskip=12mm}
\hypersetup{unicode=false, pdftoolbar=true, pdfmenubar=true, pdffitwindow=false, pdfstartview={FitH},
	pdftitle={Technische Dokumentation MCPS-Drohne},
	pdfauthor={Ingo Hjalmar Speer, Marcel Wedel, Nico Müller, Tobias Klausmann},
	pdfsubject={Technische Dokumentation MCPS-Drohne},
	pdfcreator={\LaTeX\ with package \flqq hyperref\frqq},
	pdfproducer={pdfTeX \the\pdftexversion.\pdftexrevision},
	pdfkeywords={Technische Dokumentation MCPS Drohne},
	pdfnewwindow=true,
	colorlinks=true,linkcolor=black,citecolor=black,filecolor=magenta,urlcolor=black}
\pdfinfo{/CreationDate (D:20110620133321)}
\begin{document}
\titlespacing{\section}{0pt}{12pt plus 4pt minus 2pt}{-6pt plus 2pt minus 2pt}
% Kopf- und Fusszeile
\renewcommand{\sectionmark}[1]{\markright{#1}}
\renewcommand{\leftmark}{\rightmark}
\pagestyle{fancy}
\lhead{}
\chead{}
\rhead{\thesection\space\contentsname}
\lfoot{Technische Dokumentation MCPS-Drohne}
\cfoot{}
\rfoot{ Seite \thepage}
\renewcommand{\headrulewidth}{0.4pt}
\renewcommand{\footrulewidth}{0.4pt}
% Vorspann
\renewcommand{\thesection}{\Roman{section}}
\renewcommand{\theHsection}{\Roman{section}}
\pagenumbering{roman}
% ----------------------------------------------------------------------------------------------------------
% Titelseite
% ----------------------------------------------------------------------------------------------------------
\thispagestyle{empty}
\begin{center}
	\includegraphics[width=5cm]{img/thm2.png}\\
	\vspace*{2cm}
	\Large
	\textbf{Fachbereich}\\
	\textbf{Mathematik, Naturwissenschaften und Informatik }\\
	\vspace*{2cm}
	\Huge
	\textbf{Fuzzy-Regelung Ausarbeitung}\\
	\vspace*{1.5cm}
	\textbf{Künstliche Intiligenz}\\
	\vspace*{2cm}
	
	\vfill
	\normalsize
	\newcolumntype{x}[1]{>{\raggedleft\arraybackslash\hspace{0pt}}p{#1}}
	\begin{tabular}{x{6cm}p{7.5cm}}
		\rule{0mm}{5ex}\textbf{Autoren:} 
		\newline 
		\newline Joel Bartelheimer
		\newline joel.bartelheimer@mni.thm.de
		\newline 
		\newline Nico Müller
		\newline nico.mueller@mni.thm.de
		\newline
		\rule{0mm}{5ex}\textbf{Dozenten:} & Prof. Dr.\\ 
		\rule{0mm}{5ex}\textbf{Abgabedatum:} & 31.08.2016 \\ 
	\end{tabular} 
\end{center}
\pagebreak

% ----------------------------------------------------------------------------------------------------------
% Verzeichnisse
% ----------------------------------------------------------------------------------------------------------
% TODO Typ vor Nummer
\renewcommand{\cfttabpresnum}{Tab. }
\renewcommand{\cftfigpresnum}{Abb. }
\settowidth{\cfttabnumwidth}{Abb. 10\quad}
\settowidth{\cftfignumwidth}{Abb. 10\quad}
\titlespacing{\section}{0pt}{12pt plus 4pt minus 2pt}{2pt plus 2pt minus 2pt}
\singlespacing
\rhead{INHALTSVERZEICHNIS}
\renewcommand{\contentsname}{II Inhaltsverzeichnis}
\phantomsection
\addcontentsline{toc}{section}{\texorpdfstring{II \hspace{0.35em}Inhaltsverzeichnis}{Inhaltsverzeichnis}}
\addtocounter{section}{1}
\tableofcontents
\pagebreak
\rhead{VERZEICHNISSE}
\listoffigures
\pagebreak
\listoftables
%\pagebreak
\renewcommand{\lstlistlistingname}{Listing-Verzeichnis}
{\labelsep2cm\lstlistoflistings}
\pagebreak
% ----------------------------------------------------------------------------------------------------------
% Inhalt
% ----------------------------------------------------------------------------------------------------------
% Abstände Überschrift
\titlespacing{\section}{0pt}{12pt plus 4pt minus 2pt}{-6pt plus 2pt minus 2pt}
\titlespacing{\subsection}{0pt}{12pt plus 4pt minus 2pt}{-6pt plus 2pt minus 2pt}
\titlespacing{\subsubsection}{0pt}{12pt plus 4pt minus 2pt}{-6pt plus 2pt minus 2pt}
% Kopfzeile
\renewcommand{\sectionmark}[1]{\markright{#1}}
\renewcommand{\subsectionmark}[1]{}
\renewcommand{\subsubsectionmark}[1]{}
\lhead{Kapitel \thesection}
\rhead{\rightmark}
\onehalfspacing
\renewcommand{\thesection}{\arabic{section}}
\renewcommand{\theHsection}{\arabic{section}}
\setcounter{section}{0}
\pagenumbering{arabic}
\setcounter{page}{1}
% ----------------------------------------------------------------------------------------------------------
% Einleitung
% ----------------------------------------------------------------------------------------------------------
\section{Einleitung}

\subsection{Anwendungsgebiete}

\section{Fuzzy-Logic Grundlagen}

Obwohl diese Ausarbeitung sich auf Fuzzy-Regelung spezialisiert, ist es notwendig die Grundbegriffe der Fuzzy-Logic zu verstehen. Die theoretischen Konzepte von Fuzzy-Mengen, Fuzzy-Realtionen sowie Fuzzy-Operationen werden in diesem Kapitel wiederholt.

\subsection{Fuzzy-Logic Allgemein}

In der klassichen Mengenlehre, in der Domaine von Fuzzy-Logic auch scharfe Logig genannt, hat ein jedes Element $x_i$ eine eindeutige Zugehörigkeit zu einer Menge $X$. Dies bedeutet das jedes Element $x_i$ den Zugehörigkeitsgrad $\mu$ von genau $1$ oder $0$ hat. Diese strikte Klassifizierung ist im menschlichen Verständniss nicht intuitiv ausgeprägt. Einem Menschen fällt es einfacher eine wage Aussage wie z.B. "in etwa", "relativ groß" o.Ä. zu treffen. Mithilfe der Fuzzy-Logic wird der strikte Zugehörigkeitsgrad aufgelöst, sodass jedes Element auch nur zum Teil einer Fuzzy-Menge angehören kann (siehe Abschnitte \ref{subsection:Fuzzy-Mengen}).

\label{subsection:Fuzzy-Mengen}
\subsection{Fuzzy-Mengen}

In der Domaine der Fuzzy-Logic wird eine Menge über die Zugehörigkeitsgrade ihrere Elemente definiert. Die Zugehörigkeit eines Elementes $x$ zu einer Menge $A$ wird über die Zugerhörigkeitsfunktion $\mu_A(x)$ definiert. Wichtig ist hierbei, dass jedes Element aus der Definitionsmenge $X$ einen Zugehörigkeitsgrad im reellen Wertebereich [0,1] hat. Ebenso kann jedes Element weiteren Mengen mit weiteren Zugehörigkeiten angehören.

\subsection{Fuzzy-Relationen}

\subsection{Fuzzy-Operationen}

\section{Komponenten eines Fuzzy-Reglers}

In diesem Kapitel werden die verschiedenen Komponenten einer Fuzzy-Regelung erläutert. Ein Fuzzy-Regler besteht im allgemeinen aus einem sogennanten Fuzzifizierungs-Interface, der Wissenbasis, der Entscheidungslogik sowie dem Defuzzifizierungs-Interface. Die Komponenten und deren Zusammenhänge werden in Abbildung dargestellt.


\subsection{Fuzzifizierungs-Interface}

Bei der sogennanten Fuzzifizierung geht es darum einen analogen Eingangswert 


% ----------------------------------------------------------------------------------------------------------
% Literatur
% ----------------------------------------------------------------------------------------------------------

\end{document}